\documentclass[a4paper,11pt]{article}
\usepackage[T2A]{fontenc}
\usepackage[utf8]{inputenc}
\usepackage[russian]{babel}
\usepackage{amsmath,amssymb,mathtools}
\usepackage{icomma}
\usepackage{geometry}
\geometry{left=23mm,right=23mm,top=25mm,bottom=24mm}
\usepackage{microtype}
\usepackage{enumitem}
\usepackage{needspace}
\usepackage{fancyhdr}
\usepackage{tikz}
\usepackage{xcolor}
\usepackage{hyperref}
\hypersetup{hidelinks}
\definecolor{accent}{HTML}{007D80}
\definecolor{muted}{HTML}{51575A}
\setlength{\parindent}{0pt}
\setlength{\parskip}{5pt}
\linespread{1.10}
\setlist[itemize]{leftmargin=5mm,itemsep=1pt,topsep=3pt}
\pagestyle{fancy}
\setlength{\headheight}{25pt}
\fancyhf{}
\fancyhead[L]{\small\color{muted}ОГЭ. Вариант 17}
\fancyhead[R]{\small\color{muted}Ярослав Верещагин}
\fancyfoot[C]{\small\thepage}
\renewcommand{\headrulewidth}{0.35pt}
\renewcommand{\headrule}{\hbox to\headwidth{\color{accent}\leaders\hrule height \headrulewidth\hfill}}
\newcommand{\task}[1]{\needspace{11\baselineskip}\vspace{9pt}{\Large\color{accent}Задание №\,#1}\par\vspace{1pt}\color{black}}
\newcommand{\practice}[1]{\par\vspace{4pt}{\color{accent}\textbf{Задача на закрепление.}} #1\par}
\begin{document}
\begin{titlepage}
\thispagestyle{empty}
\centering
\vspace*{27mm}
{\large\color{accent}МАТЕМАТИКА\par}
\vspace{29mm}
{\Huge Ревью домашней работы\par}
\vspace{16mm}
{\LARGE\color{accent}ОГЭ. Вариант 17\par}
\vspace{13mm}
{\large Подробный разбор ошибок и задачи на закрепление\par}
\vfill
{\Large Ярослав Верещагин\par}
\vspace{17mm}
{\large 19 ноября 2026 г.\par}
\vspace{23mm}
{\small Лёвин Артём Александрович\par}
{\small эксклюзивно для Ярослава Верещагина\par}
\vspace{17mm}
\begin{tikzpicture}
\draw[accent,line width=.75pt] (0,0) .. controls (1.4,.54) and (2.9,-.48) .. (4.4,.0)
 .. controls (5.9,.48) and (7.2,-.50) .. (8.8,0);
\fill[accent] (4.4,0) circle (1.1pt);
\end{tikzpicture}
\end{titlepage}

\task{6}
\textbf{Что не получилось.} При вычислении произведения со степенью отрицательного числа потерялся знак: слагаемое с \( (-10)^3 \) оказалось положительным. Здесь важно различать чётную и нечётную степени отрицательного основания.

\textbf{Разбор.} В условии записано
\[
0{,}8\cdot(-10)^4+5\cdot(-10)^3-76.
\]
У четвёртой степени показатель чётный, поэтому \((-10)^4=10\,000\). У третьей степени показатель нечётный, поэтому \((-10)^3=-1\,000\). Подставляем с сохранением знаков:
\[
\begin{aligned}
0{,}8\cdot10\,000+5\cdot(-1\,000)-76
 &=8\,000-5\,000-76\\
 &=\boxed{2\,924}.
\end{aligned}
\]
\textbf{Правило.} Отрицательное число в чётной степени положительно, а в нечётной остаётся отрицательным. Сначала вычисляйте степень, затем умножение и только потом сложение и вычитание.

\practice{Вычислите \(0{,}6\cdot(-10)^4+4\cdot(-10)^3-90\).}

\task{7}
\textbf{Что не получилось.} При переводе отрицательных дробей в десятичные приближения пропали минусы. Из-за этого было выбрано число \(-0{,}1\), которое не лежит между заданными границами.

\textbf{Разбор.} Нужно найти число между \(-\dfrac{2}{13}\) и \(-\dfrac{4}{15}\). Оценим обе границы:
\[
-\frac{2}{13}\approx-0{,}154,\qquad
-\frac{4}{15}\approx-0{,}267.
\]
На числовой прямой более отрицательное число находится левее. Поэтому ищем число из промежутка
\[
-0{,}267< x < -0{,}154.
\]
Подходит \(-0{,}2\), поскольку
\[
-0{,}267<-0{,}2<-0{,}154.
\]
\textbf{Верный ответ:} \(\boxed{2}\) (вариант \(-0{,}2\)).

\textbf{Правило.} При сравнении отрицательных десятичных дробей число с большим модулем меньше. Записывайте знак минус и в приближении, и на числовой прямой.

\practice{Из чисел \(-0{,}1\), \(-0{,}2\), \(-0{,}3\), \(-0{,}4\) выберите число между \(-\dfrac{3}{11}\) и \(-\dfrac{1}{6}\).}

\task{10}
\textbf{Что не получилось.} В подсчёте докладов, оставшихся на третий и четвёртый дни, неверно использованы данные первых двух дней. В решении получилось по 18 докладов, хотя по условию в первый день было 20, а во второй 12.

\textbf{Разбор.} Всего запланировано 60 докладов. Сначала вычтем доклады первых двух дней:
\[
60-20-12=28.
\]
Оставшиеся 28 докладов поровну распределены между третьим и четвёртым днями:
\[
28:2=14.
\]
Порядок докладов случаен, значит каждый из 60 докладов с одинаковой вероятностью может оказаться докладом профессора М. Благоприятных мест для последнего дня --- 14. Следовательно,
\[
P=\frac{14}{60}=\boxed{\frac{7}{30}}.
\]
\textbf{Правило.} В задачах на вероятность сначала определите общее число равновозможных исходов, затем число подходящих исходов. Прежде чем делить остаток по дням, вычтите все уже распределённые доклады.

\practice{На конференции 80 докладов. В первый день запланировано 26, во второй --- 18, остальные поровну распределены между третьим и четвёртым днями. Найдите вероятность того, что доклад случайно выбранного участника поставят на четвёртый день.}

\task{11}
\textbf{Что не получилось.} Формулы для двух парабол с ветвями вниз перепутаны: у них одинаковая высота вершины, но разные положения относительно оси \(y\).

\textbf{Разбор.} Для параболы \(y=ax^2+bx+c\) абсцисса вершины равна
\[
x_0=-\frac{b}{2a}.
\]
Проверим все три формулы по направлению ветвей и положению вершины.
\begin{itemize}
\item \(y=-2x^2+6x-2\): ветви вниз, \(x_0=-6/(-4)=1{,}5\). Вершина справа; график \textbf{Б}.
\item \(y=-2x^2-6x-2\): ветви вниз, \(x_0=6/(-4)=-1{,}5\). Вершина слева; график \textbf{В}.
\item \(y=2x^2+6x+2\): ветви вверх, \(x_0=-6/4=-1{,}5\). Это график \textbf{А}.
\end{itemize}
В порядке \(\text{А},\text{Б},\text{В}\) получаем
\[
\boxed{3\;1\;2}.
\]
\textbf{Правило.} Знак коэффициента при \(x^2\) определяет направление ветвей; формула \(x_0=-b/(2a)\) помогает отличать параболы с одинаковым направлением ветвей.

\practice{Сопоставьте функции \(1)\ y=-x^2+4x-3\), \(2)\ y=-x^2-4x-3\), \(3)\ y=x^2+4x+3\) с описаниями: А --- ветви вверх, вершина \((-2;-1)\); Б --- ветви вниз, вершина \((2;1)\); В --- ветви вниз, вершина \((-2;1)\). Запишите три номера в порядке А, Б, В.}

\task{20}
\textbf{Что не получилось.} При преобразовании уравнения \(x^6=(8-7x)^3\) возникла попытка разнести куб по разности: представить \((8-7x)^3\) как разность кубов отдельных слагаемых. Такое преобразование неверно.

\textbf{Разбор.} Заметим, что \(x^6=(x^2)^3\). Тогда
\[
(x^2)^3=(8-7x)^3.
\]
Кубическая функция строго возрастает: если кубы двух действительных чисел равны, то равны сами числа. Поэтому можно сразу приравнять основания кубов:
\[
\begin{aligned}
x^2&=8-7x,\\
x^2+7x-8&=0,\\
(x+8)(x-1)&=0.
\end{aligned}
\]
Отсюда
\[
\boxed{x=-8\quad\text{или}\quad x=1}.
\]
Проверка: при \(x=-8\) обе части равны \(64^3\); при \(x=1\) обе равны \(1\).

\textbf{Правило.} Не существует равенства \((a-b)^3=a^3-b^3\) в общем случае. Если две третьи степени равны, удобнее сравнить их основания, не раскрывая кубы.

\practice{Решите уравнение \(x^6=(5-4x)^3\).}

\vfill
\begin{center}
{\small\color{muted}Лёвин Артём Александрович\par
эксклюзивно для Ярослава Верещагина}
\end{center}
\end{document}
